\section{Experiments}
\label{sec:experiments}

\outlineblock{Evidence order}{
The section proceeds from overall efficacy to out-of-surface-form behavior,
role identification, preservation, and visible failures. This order answers
whether the method works, what explanation survives the controls, and where
the conclusion stops.}

\subsection{Experimental Setup}
\label{subsec:experimental-setup}

\outlineblock{Seven mechanisms}{
List the seven equal-status mechanisms and the 24 causal plus 18 specificity
cases per mechanism. Report 294 semantic cases, eight main streams, 2352 main
videos, and 96 additional Water/Fracture identification videos.}

\outlineblock{Models and baselines}{
Separate the Wan block (Original, \matchedcontrol, \methodshort) from the
CogVideoX block (Original, \negativeprompt, \videoeraser,
\tvtwounlearning, \safree). State why raw scores across backbones are
descriptive and why external headline contrasts require shared Original
capability.}

\outlineblock{Evaluation and statistics}{
Give the compact definitions of Original eligibility, usability, \ces, and
\su. State full-49-frame blinded review, final human-canonical authority,
case-level inference, mechanism-first equal weighting, paired bootstrap,
Holm correction, and registered non-inferiority margins. Put the full rubric,
audit flow, and formulas in \appref{app:evaluation-metrics}.}

\subsection{Main Results across Seven Mechanisms}
\label{subsec:main-results}

\outlineblock{Table 1 - seven-mechanism main view}{
Present eight methods by seven mechanisms plus macro \ces, visibly separated
into Wan and CogVideoX blocks. Report observed ranking, mechanism
heterogeneity, and the descriptive cross-backbone boundary.}

\outlineblock{Table 2 - Wan main contrast and guardrails}{
Report the single-factor \methodshort-minus-Matched contrast together with
\su, Original capability, receiver preservation, video quality, usability,
confidence intervals, multiplicity status, and registered preservation
margins. State external shared-capability estimability separately.}

\claimboundary{Use \emph{highest observed} rather than a significance claim
unless the final interval, multiplicity correction, and estimability gates all
pass. Do not convert a non-estimable external contrast into a win or a loss.}

\subsection{Role or Word? Identification Controls}
\label{subsec:role-or-word}

\outlineblock{Table 3 - two-mechanism identification design}{
Compare \methodshort\ with Generic Paraphrase and frequency-matched Bystander
Token controls. For each, report causal \ces, the implicit-footprint subset,
and specificity difference against its registered non-inferiority margin.
Explain what lexical explanation each control removes and state explicitly
that these runs cover only Water and Fracture.}

\claimboundary{If any final identification guardrail fails, describe the
evidence as partial support or as behavior consistent with role conditioning.
Do not claim that the full role-conditioned novelty criterion passed.}

\subsection{Beyond Explicit Footprints and Seen Sources}
\label{subsec:generalization}

\outlineblock{Table 4 - two generalization questions}{
Test whether the gain remains when the prompt does not name the footprint and
when the source identity is excluded from training. Report the paired
contrasts, intervals, per-mechanism counts, and their overlap. These are
behavioral generalization tests, not independent replications.}

\subsection{Preservation and Failure Analysis}
\label{subsec:preservation-failures}

\outlineblock{Continuous efficacy components}{
Separate the gated source and footprint contributions to \ces. Distinguish a
graded reduction in visibility from complete absence and strict success.}

\outlineblock{Guardrails and difficult mechanisms}{
Report \su, receiver preservation, video quality, and usable fraction against
their prespecified margins. Then discuss negative or weak mechanisms,
Original-capability limits, residual sources or footprints, and any
efficacy-preservation trade-off. Do not hide these observations in the
appendix.}

\subsection{Qualitative Comparisons}
\label{subsec:qualitative-main}

\outlineblock{Figure 3 - compact temporal evidence}{
Use identical frame indices and method order for one representative success,
one implicit-footprint case, and one failure. Label source, footprint, and
receiver. Use a deterministic metadata-only selection rule that never reads
per-method media or scores, and disclose that this rule is frozen after
aggregate unblinding. Treat the display as descriptive; move the full
eight-stream, seven-mechanism grids to
\appref{app:qualitative-failures}.}
